%% Chapter 10, Section 10.1 — Variational Inference
%% Source: A First Course in Monte Carlo Methods, Sanz-Alonso & Al-Ghattas
%% Book pages 113–118

\section{Variational Inference}
\label{sec:10.1}

The goal of Variational Inference (VI) is to find a tractable distribution $g^*$ that is close to a
given target distribution $f$, so that quantities of interest with respect to the target can be cheaply
computed using the variational distribution $g^*$. For instance, once $g^*$ is found, expected values
with respect to the target can be approximated as follows:
\[
  \calI_f[h] = \int h(x) f(x)\,dx \approx \int h(x) g^*(x)\,dx = \calI_{g^*}[h],
\]
which is useful if computing $\calI_{g^*}[h]$ is cheaper than computing $\calI_f[h]$.

The variational distribution $g^*$ is defined as the (numerical) solution to an optimization problem.
Precisely, one specifies a family $\mathcal{D}$ of tractable distributions and sets
\begin{equation}
  g^* = \argmin_{g \in \mathcal{D}}\, d_{\mathrm{KL}}(g \| f),
  \label{eq:10.1}
\end{equation}
where $d_{\mathrm{KL}}$ is the Kullback-Leibler (KL) divergence between p.d.f.s. Thus, $g^*$ is the p.d.f.\ in
$\mathcal{D}$ which is closest to the target in KL-divergence, as represented in Figure~\ref{fig:10.1}.
The family $\mathcal{D}$ of admissible distributions needs to be carefully chosen, since it determines
the accuracy and computational cost of the methodology. We will introduce the KL-divergence in
Subsection~\ref{sec:10.1.1}, and we will discuss in Subsection~\ref{sec:10.1.2} the choice of the
family $\mathcal{D}$ along with a Coordinate Ascent Variational Inference (CAVI) algorithm to solve
the optimization problem~\eqref{eq:10.1}.

\begin{figure}[htbp]
  \centering
  \includegraphics[width=0.55\textwidth]{figures/fig_10_1}
  \caption{\textit{Illustration of the variational distribution $g^*$.}}
  \label{fig:10.1}
\end{figure}

%% -------------------------------------------------------
\subsection{Kullback-Leibler Divergence and Evidence Lower-Bound}
\label{sec:10.1.1}
%% -------------------------------------------------------

In this section, we motivate the choice of the KL-divergence to quantify closeness between p.d.f.s
in~\eqref{eq:10.1}. Recall that the KL-divergence between p.d.f.s $g$ and $f$ is given by
\[
  d_{\mathrm{KL}}(g \| f) = \mathbb{E}_g\!\left[\log\!\left(\frac{g}{f}\right)\right],
\]
provided that $g$ is absolutely continuous with respect to $f$, and otherwise $d_{\mathrm{KL}}(g\|f) = \infty$.
Note that the KL-divergence satisfies $d_{\mathrm{KL}}(g\|f) \geq 0$ with equality iff $g = f$. However,
the KL-divergence is not symmetric, does not satisfy the triangle inequality, and may take value
infinity. Therefore, it is not a distance in the space of p.d.f.s.

An important motivation to work with the KL-divergence is that the optimization problem of finding
the p.d.f.\ that minimizes the KL-divergence with the target can be reformulated into an equivalent
optimization problem which does not require the target to be normalized. Indeed, we will show that
minimizing KL-divergence is equivalent to maximizing the Evidence Lower-Bound (ELBO), which we now
introduce.

\begin{definition}[Evidence Lower-Bound]
  \label{def:10.1}
  \textit{The} ELBO \textit{of a p.d.f.\ $g$ with respect to an unnormalized p.d.f.\ $\tilde{f}$
  is given by}
  \[
    \ELBO(g) \coloneqq \mathbb{E}_g[\log \tilde{f}] - \mathbb{E}_g[\log g]. \qquad \square
  \]
\end{definition}

The next result shows that the KL-divergence is equal, up to an additive constant, to the negative
ELBO\@. As a consequence, minimizing the KL-divergence is equivalent to maximizing the ELBO\@.

\begin{theorem}[Relationship Between KL-Divergence and ELBO]
  \label{thm:10.2}
  \textit{Let $f(x) = \tilde{f}(x)/c$ be a p.d.f.\ and let $\ELBO(g)$ denote the ELBO of a p.d.f.\
  $g$ with respect to the unnormalized p.d.f.\ $\tilde{f}$. Then,}
  \[
    d_{\mathrm{KL}}(g \| f) = -\ELBO(g) + \log c.
  \]
\end{theorem}

\begin{proof}
By direct calculation:
\begin{align*}
  d_{\mathrm{KL}}(g \| f)
    &= \mathbb{E}_g\!\left[\log\!\left(\frac{g}{f}\right)\right] \\
    &= \mathbb{E}_g[\log g] - \mathbb{E}_g[\log f] \\
    &= \mathbb{E}_g[\log g] - \mathbb{E}_g\!\left[\log\!\left(\frac{\tilde{f}}{c}\right)\right] \\
    &= \mathbb{E}_g[\log g] - \mathbb{E}_g\!\left[\log \tilde{f}\right] + \mathbb{E}_g[\log c] \\
    &= -\ELBO(g) + \log c. \qedhere
\end{align*}
\end{proof}

Theorem~\ref{thm:10.2} implies that the optimization problem~\eqref{eq:10.1} can be reformulated in
terms of maximizing the ELBO\@. The following corollary shows that the ELBO provides a lower-bound on
the (log-)evidence, which motivates its name. We will discuss the interpretation of the normalizing
constant $c$ as model evidence in Example~\ref{ex:10.4} below.

\begin{corollary}[Lower Bound Property of ELBO]
  \label{cor:10.3}
  \textit{In the setting of Theorem~\ref{thm:10.2}, it holds that}
  \[
    \ELBO(g) \leq \log c.
  \]
\end{corollary}

\begin{proof}
Since the KL-divergence is non-negative and $d_{\mathrm{KL}}(g\|f) + \ELBO(g) = \log c$, the result
follows.
\end{proof}

\begin{example}[VI in Bayesian Statistics]
  \label{ex:10.4}
  In Bayesian statistics, inference over a parameter $\theta$ given data $y$ is based on the
  posterior distribution
  \[
    f(\theta | y) = \frac{1}{f(y)}\, f(y|\theta)\, f(\theta),
  \]
  which combines the likelihood $f(y|\theta)$ with a prior distribution $f(\theta)$. Computing
  posterior expectations can be expensive when the dimension of $\theta$ is large. In such a case,
  computing the normalizing constant $f(y)$, called the evidence, can also be challenging, and it is
  natural to view $\tilde{f}(\theta) \coloneqq f(y|\theta)\,f(\theta)$ as an unnormalized target
  distribution with unknown normalizing constant $c \coloneqq f(y)$. Corollary~\ref{cor:10.3} then
  shows that $\ELBO(g) \leq \log f(y)$. This inequality explains the name \textit{evidence
  lower-bound}: $\ELBO(g)$ gives a lower bound on the log-evidence $\log f(y)$. In addition, this
  inequality motivates a crude but cheap model selection criteria: since a large evidence suggests
  good model fit, one may choose the model that gives the highest ELBO\@.

  In this Bayesian context, it is insightful to rewrite the ELBO as follows
  \begin{align*}
    \ELBO(g)
      &= \mathbb{E}_g[\log f(y|\theta)] + \mathbb{E}_g[\log f(\theta)] - \mathbb{E}_g[\log g(\theta)] \\
      &= \mathbb{E}_g[\log f(y|\theta)] + \mathbb{E}_g\!\left[\log\!\left(\frac{f(\theta)}{g(\theta)}\right)\right] \\
      &= \mathbb{E}_g[\log f(y|\theta)] - d_{\mathrm{KL}}(g(\theta)\|f(\theta)).
  \end{align*}
  The first term in the last displayed equation is the expected log-likelihood, whereas the second is
  the negative KL-divergence between the prior and the variational density $g$. From the above
  formulation, we obtain an intuition that the optimal $g$ should find a compromise between
  maximizing the expected log-likelihood and being close to the prior distribution. The first term
  promotes matching the data, and the second term promotes matching prior beliefs.
\end{example}

%% -------------------------------------------------------
\subsection{Mean-Field Approximation and Coordinate Ascent Variational Inference}
\label{sec:10.1.2}
%% -------------------------------------------------------

This section discusses the choice of the family $\mathcal{D}$ of admissible distributions in
\eqref{eq:10.1}. Ideally, the family $\mathcal{D}$ should be chosen so that:
\begin{enumerate}[(i)]
  \item there is $q \in \mathcal{D}$ that accurately approximates the target;
  \item expectations with the variational distribution $q^*$ can be computed efficiently. Thus,
        $\mathcal{D}$ should contain tractable distributions;
  \item the optimization problem~\eqref{eq:10.1} can be solved efficiently.
\end{enumerate}

In practice, a compromise needs to be made: enlarging the family $\mathcal{D}$ potentially helps in
the approximation accuracy, but it typically leads to a harder optimization problem. A popular choice
of $\mathcal{D}$ is the mean-field family, which sacrifices the ability to accurately represent target
distributions with highly correlated coordinates, but leads to efficient optimization and inference
algorithms.

\begin{definition}[Mean-Field Family]
  \label{def:10.5}
  \textit{A p.d.f.\ $g(x)$ with $x = (x_1, \ldots, x_d) \in \R^d$ is said to belong to the
  mean-field family if its marginals are independent, that is, if it can be written in the form}
  \[
    g(x) = \prod_{i=1}^{d} g_i(x_i),
  \]
  \textit{where $g_i$ is a one-dimensional p.d.f.\ over the $i$-th coordinate $x_i$.}
\end{definition}

In other words, the mean-field assumption requires that the variational density can be factorized as
a product of distributions. We work with the case where $g$ factorizes as a product of
one-dimensional distributions, but it is also possible to consider more general cases where it
factorizes into a product of larger groups of coordinates.

\medskip
\noindent\fbox{\parbox{\dimexpr\linewidth-2\fboxsep-2\fboxrule\relax}{%
  \textbf{Pros and Cons}: As we shall see, solving the optimization problem~\eqref{eq:10.1} is
  relatively straightforward when using the mean-field family. While restrictive, the mean-field
  family still affords some flexibility as the marginals $g_i$ are left unconstrained. If required,
  we can impose a specific parametric family on these marginals as well. The main caveat of the
  mean-field family is that it cannot capture the dependence structure of the different variables,
  and estimates computed with a mean-field distribution will be far off when the target has highly
  correlated variables. In that case, marginal variances can be severely over or under-estimated.
}}
\medskip

\begin{remark}
  \label{rem:10.6}
  \textit{The mean-field assumption has its roots in statistical physics, where `mean-field' refers
  to simplifying difficult high-dimensional stochastic models by considering simpler ones which ignore
  second-order effects. While we focus on the mean-field variational family due to its popularity and
  relative simplicity, it is possible to go beyond this restrictive assumption by incorporating
  dependencies between the marginals or by adding new latent variables. Both of these approaches
  reduce the bias of VI, but can substantially increase the difficulty of the associated optimization
  problem~\eqref{eq:10.1}.}
\end{remark}

Under the mean-field approximation, each $g_i$ can be optimized individually, and this approach is
known as Coordinate Ascent Variational Inference (CAVI). We next prove a theorem that will lead to
this technique. Let us set some notation. We denote by $x_i$ the $i$-th component of
$x \in \R^d$. We denote by $x_{-i} \in \R^{d-1}$ the vector obtained by removing the $i$-th
component of $x$. To highlight the dependence of $f(x)$ on $x_i$ and $x_{-i}$, we write
$f(x_i, x_{-i})$ with abuse of notation. Finally, $g_{-i}$ will denote a p.d.f.\ over the variables
$x_{-i}$.

\begin{theorem}[Coordinate-Wise ELBO Maximization]
  \label{thm:10.7}
  \textit{Let $g(x) = \prod_{i=1}^d g_i(x_i)$ and let $g_{-i}(x_{-i}) = \prod_{j \neq i} g_j(x_j)$
  for fixed $g_j$, $j \neq i$. Then, the $g_i$ that maximizes $\ELBO(g)$ is given by}
  \begin{equation}
    g_i^*(x_i) \propto \exp\!\Bigl(\mathbb{E}_{g_{-i}}\bigl[\log f(x_i, x_{-i})\bigr]\Bigr).
    \label{eq:10.2}
  \end{equation}
\end{theorem}

\begin{proof}
Throughout this proof, we denote by $C$ a constant which does not depend on $g_i$ and may change
from line to line. We will show that
\[
  \ELBO(g) = -d_{\mathrm{KL}}(g_i \| g_i^*) + C,
\]
where $g_i^*$ is given by~\eqref{eq:10.2}. This implies the desired result, since the $g_i$ that
minimizes the KL-divergence in the right-hand side is clearly $g_i = g_i^*$.

First, we show that
\[
  \ELBO(g) = \mathbb{E}_{g_i}\!\Bigl[\mathbb{E}_{g_{-i}}\bigl[\log f(x_i, x_{-i})\bigr]
             - \log g_i(x_i)\Bigr] + C.
\]
Indeed,
\begin{align*}
  \ELBO(g)
    &= \mathbb{E}_g[\log f(x)] - \mathbb{E}_g[\log g(x)] \\
    &= \mathbb{E}_g[\log f(x_i, x_{-i})] - \sum_i \mathbb{E}_{g_i}[\log g_i(x_i)] \\
    &= \mathbb{E}_{g_i}\!\Bigl[\mathbb{E}_{g_{-i}}\bigl[\log f(x_i, x_{-i}) \mid x_i\bigr]\Bigr]
       - \mathbb{E}_{g_i}[\log g_i(x_i)] + C,
\end{align*}
and the first term in the right-hand side can be simplified using that
\begin{align*}
  \mathbb{E}_{g_{-i}}\bigl[\log f(x_i, x_{-i}) \mid x_i\bigr]
    &= \int_{\R^{d-1}} \log f(x_i, x_{-i})\cdot g(x_{-i}|x_i)\,dx_{-i} \\
    &= \int_{\R^{d-1}} \log f(x_i, x_{-i})\cdot g_{-i}(x_{-i})\,dx_{-i} \\
    &= \mathbb{E}_{g_{-i}}\bigl[\log f(x_i, x_{-i})\bigr].
\end{align*}
Therefore,
\begin{align*}
  \ELBO(g)
    &= \mathbb{E}_{g_i}\!\Bigl[\mathbb{E}_{g_{-i}}\bigl[\log f(x_i, x_{-i})\bigr] - \log g_i(x_i)\Bigr] + C \\
    &= \mathbb{E}_{g_i}\!\Bigl[\log\!\Bigl(\exp\!\bigl(\mathbb{E}_{g_{-i}}\bigl[\log f(x_i, x_{-i})\bigr]\bigr)\Bigr)
       - \log g_i(x_i)\Bigr] + C \\
    &= -\mathbb{E}_{g_i}\!\left[\log\!\left(\frac{g_i(x_i)}{\exp\!\bigl(\mathbb{E}_{g_{-i}}[\log f(x_i,x_{-i})]\bigr)}\right)\right] + C \\
    &= -d_{\mathrm{KL}}(g_i \| g_i^*) + C,
\end{align*}
and the proof is complete.
\end{proof}

Theorem~\ref{thm:10.7} gives us an expression for the optimizer of the ELBO with respect to $g_i$
holding the other coordinates fixed. This procedure can be repeated iteratively, optimizing each
coordinate and holding the others fixed. The resulting algorithm is known as Coordinate Ascent
Variational Inference (CAVI):

\begin{algorithm}[H]
  \caption{Coordinate Ascent Variational Inference (CAVI)}
  \label{alg:10.1}
  \begin{algorithmic}[1]
    \Require Unnormalized target density $\tilde{f}$.
    \State \textbf{Initialization}: Choose initial variational densities $g_i$ for each $i = 1, \ldots, d$.
    \While{the ELBO/KL-divergence has not converged}
      \For{$i = 1, \ldots, d$}
        \State Update
        \[
          g_i(x_i) \propto \exp\!\Bigl(\mathbb{E}_{g_{-i}}\bigl[\log \tilde{f}(x_i, x_{-i})\bigr]\Bigr).
        \]
      \EndFor
    \EndWhile
    \Ensure $g(x) = \prod_{i=1}^{d} g_i(x_i) \approx f(x)$.
  \end{algorithmic}
\end{algorithm}

In each update, CAVI lowers the KL-divergence between $q^*$ and the target $f$, and hence one can
expect convergence to a local minimizer under suitable assumptions. The ELBO is not necessarily
concave and therefore there is in general no guarantee to reach a maximizer of the ELBO, or,
equivalently, a minimizer of the KL-divergence.

\medskip
\noindent\fbox{\parbox{\dimexpr\linewidth-2\fboxsep-2\fboxrule\relax}{%
  \textbf{Pros and Cons}: Each CAVI update lowers the KL-divergence, which facilitates introducing
  stopping criteria for the algorithm. If the factors $g_i$ are assumed to be in the exponential
  family, computing each factor update is simplified significantly. However, since CAVI relies on the
  mean-field assumption, it inherits all its disadvantages. In addition, for some target distributions
  it may be computationally expensive to compute the updates of the variational factors. In such a
  case, other optimization methods (e.g.\ stochastic optimization and natural gradient methods) rather
  than CAVI can be used to maximize the ELBO\@.
}}
\medskip

\begin{remark}
  \label{rem:10.8}
  \textit{CAVI shares a similar structure with the Gibbs sampler. To see why, notice that}
  \[
    g_i^*(x_i) \propto \exp\!\Bigl(\mathbb{E}_{g_{-i}}\bigl[\log f(x_i, x_{-i})\bigr]\Bigr)
    \propto \exp\!\Bigl(\mathbb{E}_{g_{-i}}\bigl[\log f(x_i | x_{-i})\bigr]\Bigr).
  \]
  \textit{In Gibbs sampling, we update each coordinate by sampling the corresponding full
  conditional. CAVI uses full conditionals as well, with each density $g_i$ set to be proportional
  to the exponential of the expectation of the log of the full conditional.}
\end{remark}
